\documentclass[11pt]{article}
\usepackage[margin=1in]{geometry}
\usepackage{amsmath,amssymb,amsthm,mathtools}
\usepackage{booktabs,array}
\usepackage{microtype}
\usepackage{setspace}
\usepackage{natbib}
\usepackage[hidelinks]{hyperref}
\usepackage{enumitem}
\usepackage{graphicx}
\usepackage{xcolor}

\onehalfspacing

\newtheorem{proposition}{Proposition}
\newtheorem{theorem}{Theorem}
\newtheorem{corollary}{Corollary}
\newtheorem{lemma}{Lemma}
\newtheorem{remark}{Remark}

\newcommand{\E}{\mathbb{E}}
\newcommand{\R}{\mathbb{R}}
\newcommand{\1}{\mathbf{1}}
\newcommand{\PN}{A_N^{N}}
\newcommand{\PP}{A_N^{P}}

\title{When Supplier Competition Worsens Policy Coordination:\\Procurement Rents and Regional Industrial-Policy Specialization}
\author{Anonymous manuscript}
\date{September 2026}

\begin{document}
\maketitle

\begin{abstract}
Supplier competition is usually valued because it lowers procurement costs and raises the surplus generated by a project. We show that the same competitive force can worsen decentralized industrial-policy coordination. Two regional governments choose between a locally attractive activity and a complementary supplier-side activity. If the regions choose complementary roles, suppliers compete in a procurement auction. More suppliers lower the expected minimum production cost but, under a standard decreasing-reversed-hazard-rate condition, also dissipate the information rent accruing to the winning supplier. When the complementary region captures part of this supplier rent, stronger competition therefore has opposite effects on coordinated and decentralized policy incentives. It lowers the project scale at which a coordinated planner prefers regional differentiation, but raises the scale required for the complementary region to accept that role voluntarily. The interval with unique decentralized priority duplication consequently expands with supplier competition. Moreover, sufficiently strong competition can eliminate differentiated Nash equilibria even when regional differentiation remains socially optimal. The procurement rent-dissipation result is standard; the contribution is to show how it changes the equilibrium set of decentralized industrial policy.
\end{abstract}

\noindent\textbf{Keywords:} industrial policy; procurement; supplier competition; regional policy; coordination; auctions.\\
\textbf{JEL:} D44; H57; L13; L52; R58.

\section{Introduction}
\label{sec:introduction}

Competition among suppliers is a standard prescription in public procurement. A larger set of bidders tends to lower the minimum production cost, reduce procurement payments, and improve productive efficiency conditional on a project being undertaken. Those gains are central to auction theory and to procurement policy. Yet procurement competition also changes the incidence of surplus. In particular, stronger competition can compress the information rent earned by the winning supplier. This paper asks whether that incidence effect can feed back into a different policy margin: whether a regional government is willing to undertake the supplier-side activity required for an interregional industrial project.

The question matters whenever industrial policy is organized around complementary roles rather than a single undifferentiated subsidy. Governments increasingly use procurement, local-content initiatives, supply-chain policies, and targeted industrial support to build or localize complementary stages of production. Public procurement is itself recognized as an industrial-policy instrument, including in innovation policy \citep{ChiappinelliGiuffridaSpagnolo2025}. The OECD's recent review documents procurement measures used to support domestic and local production and emphasizes the tension among industrial-policy objectives, market access, competition, and supply-chain resilience \citep{OECD2026}. Territorial procurement is also empirically relevant: local and devolved governments differ systematically in the extent to which their procurement is sourced within their own regions \citep{EckersleyEtAl2023}. Recent industrial-organization work on industrial policy likewise emphasizes that policy effects depend on market structure, that localizing supply chains can be central, and that government policy may have a coordination role across market actors \citep{VanBiesebroeckZhang2026,MottaPolo2026}. In such settings, a policy environment with more supplier competition can improve production efficiency while changing the local rents that make a jurisdiction willing to host the supplier-side activity.

We study the smallest model that isolates this interaction. Two regional governments simultaneously choose between a locally attractive activity, denoted $U$, and a complementary supplier-side activity, denoted $D$. The direct local return to $U$ exceeds that to $D$ by $\Delta>0$. If the regions choose complementary roles, $(U,D)$ or $(D,U)$, the $U$ region is the lead or buyer side of the complementary project and the $D$ region hosts or develops the supplier-side industrial base. The $D$ side then hosts a procurement market with $N\geq2$ symmetric suppliers. Suppliers have independent private costs and compete for a project whose value is $v$. The baseline uses a second-price reverse auction because it makes the surplus decomposition transparent: expected total match surplus depends on the lowest cost, buyer surplus depends on the second-lowest cost, and winning supplier rent is the gap between the second- and first-order statistics.

The lower-stage comparative static is deliberately standard. Let $C_{1:N}$ and $C_{2:N}$ denote the lowest and second-lowest supplier costs. Define
\begin{equation}
G_N\equiv v-\E[C_{1:N}],
\qquad
R_N\equiv \E[C_{2:N}-C_{1:N}].
\label{eq:intro-GR}
\end{equation}
The first object is expected real surplus from an active complementary project and the second is expected winning supplier information rent. More suppliers reduce the expected minimum cost, so $G_N$ rises. Under the standard decreasing-reversed-hazard-rate (DRHR) condition used in procurement theory, expected winning rent $R_N$ strictly falls with the number of suppliers and converges to zero \citep{Li2005,XuLi2019}; related bidder-attraction results appear in \citet{Szech2011}. We do not claim this rent-dissipation result as new.

The new result begins when these two effects enter the regional policy game. Suppose the $D$ region internalizes a fraction $\rho\in(0,1]$ of local supplier rent. A coordinated policymaker values the full real surplus $G_N$ created by regional differentiation. The $D$ government, however, gives up its direct premium $\Delta$ and obtains only $\rho R_N$ from the supplier-side role. Let $A>0$ scale the mass or importance of the complementary project. The coordinated differentiation threshold is
\begin{equation}
A_N^P=\frac{\Delta}{G_N},
\label{eq:intro-AP}
\end{equation}
whereas the decentralized threshold at which a region is willing to accept $D$ against a region choosing $U$ is
\begin{equation}
A_N^N=\frac{\Delta}{\rho R_N}.
\label{eq:intro-AN}
\end{equation}
At a fixed $N$, this is an incidence wedge: the policymaker who values the full project switches earlier than the jurisdiction that captures only part of its benefit. That fixed-market observation is not the contribution of the present paper. It maps directly into the familiar logic that incomplete local capture delays decentralized adjustment.

Our contribution is the comparative static in market structure. Because supplier competition raises $G_N$ but lowers $R_N$, it moves the two thresholds in opposite directions:
\begin{equation}
A_{N+1}^P<A_N^P,
\qquad
A_{N+1}^N>A_N^N.
\label{eq:intro-divergence}
\end{equation}
Competition therefore makes differentiation socially desirable at a lower project scale while making the supplier-side region require a larger project before voluntarily accepting that role. The set of project scales for which decentralized policy uniquely duplicates $U$ even though coordinated policy differentiates consequently expands with supplier competition.

The stronger result concerns the equilibrium set rather than a threshold gap. For a fixed $N$, the induced government game is an anti-coordination game only when $A\rho R_N>\Delta$. In that region, $(U,D)$ and $(D,U)$ are differentiated pure Nash equilibria. Since $R_N\to0$ as supplier competition becomes sufficiently strong, those differentiated equilibria eventually disappear. At the same time, the planner's valuation of differentiation rises because $G_N$ increases. If
\begin{equation}
A(v-c_0)>\Delta,
\label{eq:intro-limit-condition}
\end{equation}
where $c_0$ is the lower endpoint of the cost support, then sufficiently large $N$ yields
\begin{equation}
\text{coordinated optimum}=\{(U,D),(D,U)\},
\qquad
\text{unique Nash equilibrium}=(U,U).
\label{eq:intro-reversal}
\end{equation}
Thus competition can destroy efficient differentiated equilibria and leave priority duplication as the unique decentralized outcome. We call this \emph{competition-induced policy reversal}.

The result is counterintuitive because the supplier market becomes more efficient conditional on the project being undertaken. The failure arises at the policy-selection layer. More competition reduces production cost but also reallocates surplus away from the local supplier side. If local governments respond to locally captured income or rents, the same increase in supplier-market competition that improves procurement performance can weaken the incentive of one region to specialize in the activity required for interregional complementarity. The policy problem is therefore not that competition directly lowers welfare. It is that competition changes the incidence supporting decentralized coordination.

The paper connects four literatures. First, auction and procurement theory studies bidder competition, participation, rent, and procurement preferences. \citet{Li2005} establishes monotonicity of expected rent under DRHR; \citet{Szech2011} studies costly bidder attraction and the difference between private and social bidder-count objectives; and \citet{XuLi2019} applies related order-statistic results to procurement competition and welfare. Endogenous participation can overturn simple competition intuitions, as \citet{LiZheng2009} show, which is why the number of suppliers is a market-structure primitive in our baseline rather than an endogenous entry choice. Procurement preferences and set-asides also redistribute surplus and alter participation \citep{Vagstad1995,Marion2007,AtheyCoeyLevin2013}. Regional procurement restrictions can explicitly trade off participation against local-industry objectives \citep{SaitoMitsuta2012}, local-content policy design can depend on local-supplier competitiveness \citep{Macatangay2016}, and actual local procurement displays a territorial sourcing dimension \citep{EckersleyEtAl2023}. Our paper takes these procurement-side forces as inputs and asks how they change a strategic public-policy game.

Second, industrial-organization research increasingly studies industrial policy as a function of market structure. \citet{AghionEtAl2015} document an interaction between industrial policy and competition, while the recent IJIO industrial-policy agenda emphasizes market structure, supply-chain localization, and coordination \citep{VanBiesebroeckZhang2026}. \citet{MottaPolo2026} provide a recent theoretical example in which market organization and appropriability shape precautionary industrial policy under supply-chain disruption. Our mechanism is different: supplier competition changes local rent incidence and thereby alters which region is willing to take a complementary policy role.

Third, place-based and local-government policy can itself be understood through competition, differentiation, and rent allocation. In an earlier regional-competition model, \citet{JustmanThisseVanYpersele2005} show that regions may differentiate infrastructure to soften competition for investors and that informational frictions can make fiscal competition inefficient. \citet{Slattery2025} models state and local governments bidding for firms and shows that intergovernmental competition can transfer rents to firms. \citet{LinLi2026} study local industrial-policy competition in a quantitative spatial equilibrium environment. We study a different direction of strategic influence: competition among suppliers changes the equilibrium of a game among governments, and can eliminate differentiated regional-policy equilibria rather than generate differentiation as a direct strategic response between regions.

Finally, the paper is related to a companion analysis of fixed-capacity regional industrial-policy composition \citep{Companion2026}. That paper takes complementary surplus and its local incidence largely as primitives and studies a continuous portfolio game. Here the policy choice is binary and the incidence wedge is generated by an explicit private-information procurement market. At any fixed number of suppliers, the unilateral switching condition can be mapped to the earlier reduced-form logic; we do not treat that fixed-market comparison as a new result. The new object is the supplier-market comparative static: stronger competition changes both total surplus and local rent capture and can alter the exact Nash-equilibrium set of regional policy.

The remainder of the paper proceeds as follows. Section \ref{sec:institutional} motivates the procurement--industrial-policy link. Section \ref{sec:model} presents the model and procurement continuation. Section \ref{sec:equilibrium} characterizes decentralized regional-policy equilibrium, and Section \ref{sec:planner} defines the coordinated fixed-instrument benchmark. Section \ref{sec:main-results} establishes competition--coordination divergence and competition-induced policy reversal. Section \ref{sec:robustness} studies incidence, trade failure, and auction-format robustness. Section \ref{sec:discussion} develops empirical implications and scope conditions. Section \ref{sec:conclusion} concludes. Proofs are in the Appendix.

\section{Institutional Motivation: Procurement as Industrial Policy}
\label{sec:institutional}

Public procurement is not only a purchasing technology. It can also determine which firms, technologies, and local production capabilities are sustained by public demand. The industrial-policy role of procurement is particularly visible in innovation-oriented procurement, local-content rules, strategic supply-chain initiatives, and government efforts to build domestic or regional production capacity. The recent survey by \citet{ChiappinelliGiuffridaSpagnolo2025} treats public procurement explicitly as an innovation-policy instrument and emphasizes design, frictions, and its interaction with the broader policy mix. The OECD likewise documents procurement measures used to support domestic and local production, employment, strategic sectors, and supply-chain resilience \citep{OECD2026}.

A recurrent policy tension is that competition and localization need not have identical incidence. In the Japanese local-government setting, \citet{SaitoMitsuta2012} examine regional eligibility requirements that promote local industry and employment while restricting outside participation. Competitive tendering can reduce expected procurement costs and select more efficient suppliers, yet it can also compress supplier rents. Procurement rules that favor local or small suppliers make this tension explicit. \citet{Vagstad1995} analyzes procurement discrimination when governments care about local firms' profits. \citet{Marion2007} and \citet{AtheyCoeyLevin2013} show empirically and structurally that procurement preferences and set-asides change entry, bidder profit, procurement cost, and efficiency. These papers study procurement design. Our paper instead holds the procurement rule fixed and asks how the competitiveness of the supplier market feeds back into a separate government decision over industrial-policy roles.

The distinction is useful for regional policy. Suppose a cross-regional project has a lead or buyer side, denoted $U$, and a complementary supplier-side industrial base, denoted $D$. A differentiated policy profile means that one jurisdiction supports the lead/buyer role while the other hosts or develops the supplier base from which procurement is conducted. If the supplier-side regional government internalizes only part of the supplier surplus generated by that procurement relationship, then the incidence of procurement surplus affects its willingness to specialize. A highly competitive procurement market may be desirable for a central buyer while reducing the local payoff that sustains decentralized specialization.

The paper therefore does not model procurement competition as a government choice. The number of suppliers $N$ is an exogenous market-structure parameter. This keeps the object distinct from the bidder-attraction problem studied by \citet{Szech2011} and \citet{XuLi2019}, and from endogenous-entry procurement models such as \citet{LiZheng2009}. The comparative static asks how different supplier-market structures change the equilibrium of the regional policy game. Endogenizing supplier entry would introduce an additional participation equilibrium and is outside the model.

This focus also matches the recent industrial-organization framing of industrial policy. \citet{VanBiesebroeckZhang2026} emphasize that industrial-policy effects depend on prevailing market structure, the localization of supply chains, and coordination across market actors. \citet{MottaPolo2026} likewise show that the appropriate industrial-policy response to supply-chain risk depends on market organization and appropriability. The present paper isolates a different channel: market competition affects the geographic incidence of rents, and that incidence changes decentralized policy specialization.

\section{Model and Procurement Continuation}
\label{sec:model}

\subsection{Regional policy choices}

There are two symmetric regions, indexed by $i\in\{1,2\}$. Each regional government chooses one of two industrial-policy roles,
\begin{equation}
s_i\in\{U,D\}.
\end{equation}
Activity $U$ is the locally attractive lead/buyer-side activity and yields direct real surplus $\Delta>0$ to the region choosing it. Activity $D$ is the complementary supplier-side activity: the region choosing it hosts or develops the supplier base used by the lead side. The direct local return to $D$ apart from procurement-side surplus is normalized to zero. The normalization only fixes the opportunity cost of taking the complementary role: a government that switches from $U$ to $D$ gives up $\Delta$.

The policy roles are technologically complementary across regions. A procurement continuation is activated only when the policy profile is differentiated, $(U,D)$ or $(D,U)$. If both governments choose $U$, each obtains its direct return $\Delta$ and no complementary project occurs. If both choose $D$, neither obtains the direct $U$ return and no complementary project occurs.

A parameter $A>0$ scales the mass or economic importance of the complementary relationship. It can be interpreted as the scale, mass, or economic importance of otherwise identical complementary procurement opportunities. It is a continuous intensity parameter, not literally a count of projects. It does not change the supplier cost distribution.

\subsection{Supplier market}

Conditional on a differentiated policy profile, the $U$ region is the lead/buyer side of the project and procures from a supplier base hosted or developed by the $D$ region. The $D$ region has $N\ge2$ potential suppliers. Supplier $k$ has private production cost
\begin{equation}
C_k\overset{iid}{\sim}F
\end{equation}
on compact support $[c_0,\bar c]$, with continuous strictly positive density $f$. Costs are independent private information and suppliers are risk neutral. The project creates gross value $v$ and the baseline assumes
\begin{equation}
v>\bar c,
\label{eq:guaranteed-trade}
\end{equation}
so trade always occurs once a differentiated policy profile is selected.

For transparency, the canonical implementation is a second-price reverse auction. The lowest-cost supplier wins and is paid the second-lowest cost. Let $C_{1:N}$ and $C_{2:N}$ denote the first and second order statistics. Truthful cost revelation is a dominant strategy, and the realized decomposition is
\begin{align}
\text{real match surplus} &= v-C_{1:N},\\
\text{buyer surplus} &= v-C_{2:N},\\
\text{winning supplier rent} &= C_{2:N}-C_{1:N}.
\end{align}
Define the expected objects
\begin{align}
m_N &\equiv \E[C_{1:N}],\label{eq:mN}\\
s_N &\equiv \E[C_{2:N}],\label{eq:sN}\\
R_N &\equiv s_N-m_N=\E[C_{2:N}-C_{1:N}],\label{eq:RN}\\
B_N &\equiv v-s_N,\label{eq:BN}\\
G_N &\equiv v-m_N=B_N+R_N.\label{eq:GN}
\end{align}
Here $G_N$ is expected real surplus of the active complementary relation, $R_N$ is expected winning supplier information rent, and $B_N$ is expected buyer surplus.

The baseline distribution satisfies decreasing reversed hazard rate (DRHR):
\begin{equation}
\frac{f(c)}{F(c)}\quad\text{is decreasing in }c.
\label{eq:drhr}
\end{equation}
This is a standard regularity condition in procurement models. We use it only for the comparative static of supplier rent. In particular, procurement theory establishes that, under DRHR, expected winning rent $R_N$ is strictly decreasing in $N$ and converges to zero \citep{Li2005,XuLi2019}. That result is imported rather than derived as a contribution of this paper.

\subsection{Regional incidence}

The $U$ region internalizes the buyer-side surplus $B_N$. The $D$ regional government internalizes a fraction
\begin{equation}
\rho\in(0,1]
\end{equation}
of winning supplier rent. Thus its expected local benefit from the procurement continuation is $\rho R_N$. The parameter $\rho$ is an incidence mapping: it is the share or weight of supplier surplus internalized by the $D$ regional government. The remaining $(1-\rho)R_N$ is outside that regional government's objective but, in the baseline, remains within the constituency of the coordinated benchmark. It is therefore not destroyed and is not interpreted mechanically as foreign ownership.

The coordinated benchmark introduced below values the full real surplus $G_N$. Thus $\rho<1$ represents incomplete regional-government internalization within the modeled supralocal constituency, not a social waste term.

\subsection{Timing and induced payoff matrix}

The timing is:
\begin{enumerate}[label=(\arabic*)]
\item $N,F,\Delta,A,v,$ and $\rho$ are common knowledge.
\item Regional governments simultaneously choose $U$ or $D$.
\item If the profile is differentiated, supplier costs are realized privately and the procurement auction is played.
\item Production and payoffs are realized.
\end{enumerate}

Solving the procurement continuation yields the expected first-stage payoff matrix in Table \ref{tab:payoffs}.

\begin{table}[ht]
\centering
\caption{Expected regional payoffs}
\label{tab:payoffs}
\begin{tabular}{c c c}
\toprule
& Region 2: $U$ & Region 2: $D$\\
\midrule
Region 1: $U$ & $(\Delta,\Delta)$ & $(\Delta+A B_N,\; A\rho R_N)$\\
Region 1: $D$ & $(A\rho R_N,\; \Delta+A B_N)$ & $(0,0)$\\
\bottomrule
\end{tabular}
\end{table}

Because $v>\bar c\ge C_{2:N}$ almost surely, $B_N>0$. Hence $U$ is always a strict best response to an opponent choosing $D$. The strategic issue is whether a government is willing to choose $D$ when the other region chooses $U$.

\section{Decentralized Regional-Policy Equilibrium}
\label{sec:equilibrium}

Against an opponent choosing $U$, a government obtains $\Delta$ from choosing $U$ and $A\rho R_N$ from choosing $D$. Define the decentralized differentiation threshold
\begin{equation}
A_N^N\equiv \frac{\Delta}{\rho R_N}.
\label{eq:AN-main}
\end{equation}
The threshold is finite because $R_N>0$ for a nondegenerate continuous cost distribution and finite $N$.

\begin{proposition}[Exact decentralized equilibrium set]
\label{prop:nash}
Fix $N\ge2$.
\begin{enumerate}[label=(\roman*)]
\item If $A<A_N^N$, $U$ is a strictly dominant action and the unique Nash equilibrium is $(U,U)$.
\item If $A=A_N^N$, let $p_i$ denote the probability that region $i$ chooses $U$. The full mixed-strategy Nash set is
\begin{equation}
\left\{(p_1,p_2)\in[0,1]^2:\ p_1=1\ \text{or}\ p_2=1\right\}.
\label{eq:boundary-nash}
\end{equation}
Its pure-strategy members are $(U,U)$, $(U,D)$, and $(D,U)$.
\item If $A>A_N^N$, the differentiated profiles $(U,D)$ and $(D,U)$ are pure Nash equilibria. In addition, there is a unique completely mixed equilibrium, which is symmetric and in which each government chooses $U$ with probability
\begin{equation}
p_N^*=\frac{\Delta+A B_N}{A(B_N+\rho R_N)}\in(0,1).
\label{eq:mixed}
\end{equation}
\end{enumerate}
\end{proposition}

The game is therefore a dominance game below the threshold and an anti-coordination game above it. At the boundary, the equilibrium correspondence contains the three pure equilibria and semi-mixed profiles in which at least one region chooses $U$ with probability one. The role of procurement is not to create an additional strategic action for the governments; it determines the continuation value $\rho R_N$ associated with accepting the $D$ role.

The mixed equilibrium provides another way to see the effect of local incidence. A larger $\rho R_N$ makes $D$ more attractive against $U$, lowering the probability with which each government chooses $U$ in the symmetric mixed equilibrium. Conversely, supplier-rent compression pushes the mixed equilibrium toward duplication. We do not use this observation as the main comparative static because the strongest result below concerns the disappearance of the differentiated pure equilibria themselves.

\begin{remark}[Fixed-$N$ interpretation]
\label{rem:fixedN}
For a fixed supplier market, the unilateral move from $(U,U)$ to $(U,D)$ is profitable for the $D$ government iff $A\rho R_N>\Delta$. This is a standard incomplete-capture switching inequality. The novel results of the paper concern how the entire equilibrium set changes when the supplier-market primitive $N$ changes.
\end{remark}

\section{Coordinated Fixed-Instrument Benchmark}
\label{sec:planner}

We compare decentralized policy with a coordinator who chooses the two regional roles while holding fixed the policy instrument set, supplier market, procurement technology, information structure, project value, and project scale. This is a coordinated fixed-instrument benchmark, not an unrestricted first best.

Procurement payments are transfers within the modeled economy. The real surplus of a differentiated project is therefore $v-C_{1:N}$, with expectation $G_N$. The aggregate modeled values of the four policy profiles are
\begin{align}
V(U,U)&=2\Delta,\label{eq:planner-UU}\\
V(U,D)=V(D,U)&=\Delta+A G_N,\label{eq:planner-UD}\\
V(D,D)&=0.\label{eq:planner-DD}
\end{align}
Define the coordinated differentiation threshold
\begin{equation}
A_N^P\equiv \frac{\Delta}{G_N}.
\label{eq:AP-main}
\end{equation}

\begin{proposition}[Exact coordinated allocation]
\label{prop:planner}
Fix $N\ge2$.
\begin{enumerate}[label=(\roman*)]
\item If $A<A_N^P$, $(U,U)$ is the unique coordinated optimum.
\item If $A=A_N^P$, $(U,U)$, $(U,D)$, and $(D,U)$ are coordinated optima.
\item If $A>A_N^P$, the exact coordinated optima are $(U,D)$ and $(D,U)$.
\end{enumerate}
\end{proposition}

Because $G_N=B_N+R_N$ and $B_N>0$, while $\rho\le1$, we have
\begin{equation}
G_N>\rho R_N.
\label{eq:capture-gap}
\end{equation}
Hence $A_N^P<A_N^N$ for every fixed $N$. This fixed-$N$ threshold ordering is not the headline result. The central question is how the two thresholds respond to supplier competition.

\section{Supplier Competition and Policy Regimes}
\label{sec:main-results}

We now vary the exogenous number of suppliers. The first comparative static is elementary: adding an independent supplier lowers the expected minimum cost strictly, so
\begin{equation}
m_{N+1}<m_N
\quad\Longrightarrow\quad
G_{N+1}>G_N.
\label{eq:G-increase}
\end{equation}
The second comparative static comes from procurement theory. Under DRHR, expected winning supplier rent is strictly decreasing and converges to zero \citep{Li2005,XuLi2019}:
\begin{equation}
R_{N+1}<R_N,
\qquad
\lim_{N\to\infty}R_N=0.
\label{eq:R-decrease}
\end{equation}
The economic contribution comes from combining these known lower-stage effects with the government game.

\subsection{Competition--coordination divergence}

\begin{theorem}[Competition--coordination divergence]
\label{thm:divergence}
Under the baseline assumptions,
\begin{equation}
A_{N+1}^P<A_N^P
\qquad\text{and}\qquad
A_{N+1}^N>A_N^N.
\label{eq:opposite-thresholds}
\end{equation}
Thus stronger supplier competition moves the coordinated and decentralized differentiation thresholds in opposite directions.
\end{theorem}

The first inequality follows because competition raises expected real project surplus. The second follows because competition dissipates the locally captured supplier rent that rewards a region for accepting $D$. The same change in market structure therefore makes differentiation more attractive to a coordinator and less attractive to the decentralized region that must specialize on the supplier side.

Define the inefficient-duplication interval
\begin{equation}
\mathcal W_N\equiv(A_N^P,A_N^N).
\label{eq:wedgeN}
\end{equation}
For every $A\in\mathcal W_N$, Proposition \ref{prop:nash} implies that $(U,U)$ is the unique Nash equilibrium while Proposition \ref{prop:planner} implies that the coordinator strictly chooses a differentiated profile.

\begin{corollary}[Strict expansion of unique inefficient duplication]
\label{cor:wedge-expansion}
Under the baseline assumptions,
\begin{equation}
\mathcal W_N\subsetneq \mathcal W_{N+1}.
\label{eq:wedge-nesting}
\end{equation}
\end{corollary}

This result differs from a conventional statement that competition improves or worsens welfare. Conditional on differentiation, procurement becomes more efficient as $N$ rises. What expands is the set of environments in which decentralized governments fail to select the differentiated configuration in the first place.

\subsection{Competition-induced policy reversal}

The threshold result has a sharper equilibrium implication. Let $c_0$ be the lower endpoint of supplier costs. Under the maintained compact support and positive density, the minimum order statistic converges to $c_0$, and hence
\begin{equation}
\lim_{N\to\infty}G_N=v-c_0.
\label{eq:Glimit}
\end{equation}
At the same time, DRHR implies $R_N\to0$.

\begin{theorem}[Competition-induced policy reversal]
\label{thm:reversal}
Suppose
\begin{equation}
A(v-c_0)>\Delta.
\label{eq:reversal-condition}
\end{equation}
Then there exists a finite $\bar N$ such that for all $N\ge\bar N$,
\begin{equation}
\operatorname{argmax}V=\{(U,D),(D,U)\},
\qquad
NE_N=\{(U,U)\}.
\label{eq:reversal-result}
\end{equation}
\end{theorem}

Theorem \ref{thm:reversal} is the main result. If the asymptotic real surplus from a differentiated project is large enough to cover the sacrificed direct premium, a coordinator eventually prefers differentiation. Yet competition drives locally captured supplier rent toward zero, so the $D$ role eventually becomes privately unattractive to either regional government. The market can therefore approach productive efficiency within the procurement subgame while decentralized policy converges toward the wrong policy configuration.

A useful corollary starts from a market in which differentiated Nash equilibria already exist.

\begin{corollary}[Destruction of differentiated Nash equilibria]
\label{cor:destruction}
Suppose that for some $N_0$,
\begin{equation}
A\rho R_{N_0}>\Delta
\label{eq:initial-diff}
\end{equation}
so that $(U,D)$ and $(D,U)$ are pure Nash equilibria. Then there exists a finite $N_1>N_0$ such that for all sufficiently large $N\ge N_1$, those differentiated equilibria no longer exist and $(U,U)$ is the unique Nash equilibrium. If the coordinator differentiates at $N_0$, it continues to differentiate for every $N>N_0$.
\end{corollary}

This corollary gives the result an equilibrium-set interpretation. Supplier competition need not merely widen an inefficiency wedge around a fixed equilibrium. It can change the qualitative regime from one with efficient differentiated equilibria to one with unique duplicated priorities.

\subsection{Uniform illustration}

For intuition, let $F$ be uniform on $[0,1]$ and set $\rho=1$. Then
\begin{equation}
m_N=\frac{1}{N+1},
\qquad
R_N=\frac{1}{N+1},
\qquad
B_N=v-\frac{2}{N+1}.
\label{eq:uniform-objects}
\end{equation}
The thresholds are
\begin{equation}
A_N^P=\frac{\Delta(N+1)}{v(N+1)-1},
\qquad
A_N^N=\Delta(N+1).
\label{eq:uniform-thresholds}
\end{equation}
Thus the decentralized threshold grows linearly with supplier count, whereas the coordinated threshold falls toward $\Delta/v$. Figure \ref{fig:uniform-regimes} illustrates the opposite movements for $\Delta=1$ and $v=2$. The horizontal line $A=5$ shows the equilibrium-set transition directly: differentiated pure equilibria exist at low $N$, the decentralized threshold is reached at $N=4$, and for $N\ge5$ the unique Nash equilibrium is duplication, while coordinated differentiation remains optimal throughout.

\begin{figure}[ht]
\centering
\includegraphics[width=0.82\textwidth]{../figures/uniform_regime_thresholds.pdf}
\caption{Competition moves the coordinated and decentralized differentiation thresholds in opposite directions. Uniform costs on $[0,1]$, $\rho=1$, $\Delta=1$, and $v=2$. The fixed line $A=5$ illustrates the destruction of differentiated Nash equilibria as $N$ rises.}
\label{fig:uniform-regimes}
\end{figure}

The uniform case is only an illustration; Theorems \ref{thm:divergence} and \ref{thm:reversal} use the general DRHR class.

\section{Robustness and Scope}
\label{sec:robustness}

The main theorem uses a deliberately transparent incidence structure and guaranteed trade. This section records which parts can be relaxed without changing the central message and which cannot.

\subsection{General regional incidence}

The baseline assumes that the $D$ government internalizes a fraction $\rho$ of supplier rent but none of the buyer surplus. To allow broader incidence, suppose instead that it also internalizes fraction $\sigma\ge0$ of buyer surplus. Its local benefit from accepting $D$ becomes
\begin{equation}
L_N=\rho R_N+\sigma B_N.
\label{eq:general-incidence}
\end{equation}
Under DRHR, define
\begin{equation}
d_m\equiv m_N-m_{N+1}>0,
\qquad
d_R\equiv R_N-R_{N+1}>0.
\end{equation}
Since $B_{N+1}-B_N=d_m+d_R$, local captured surplus falls when competition increases iff
\begin{equation}
\rho>\sigma\left(1+\frac{d_m}{d_R}\right).
\label{eq:incidence-condition}
\end{equation}
Thus complete local ownership of suppliers is not required. What matters is whether supplier-rent exposure is sufficiently strong relative to the region's exposure to buyer-side gains from competition.

For $F=U[0,1]$, $d_m=d_R$, so condition \eqref{eq:incidence-condition} simplifies to
\begin{equation}
\rho>2\sigma.
\label{eq:uniform-incidence}
\end{equation}
The baseline $\sigma=0$, $\rho>0$ satisfies the condition automatically.

\subsection{Trade failure and a reserve value}

The guaranteed-trade assumption $v>\bar c$ is substantive for the strong one-step monotonicity result. If $v$ lies inside the support, procurement occurs only when the lowest cost is below $v$. Define
\begin{align}
G_N^R&\equiv \E[(v-C_{1:N})_+],\label{eq:reserve-G}\\
R_N^R&\equiv \E\left[(\min\{C_{2:N},v\}-C_{1:N})\1\{C_{1:N}<v\}\right].\label{eq:reserve-R}
\end{align}
Adding suppliers raises the probability that trade occurs and can therefore increase expected supplier rent at low $N$, even though conditional competition compresses rents. Stepwise monotonicity of $R_N^R$ is not general. For example, with uniform costs on $[0,1]$ and $v=0.1$, expected supplier rent initially rises with $N$.

The asymptotic reversal is more robust. Under the same compact-support assumptions and $v>c_0$, the first two order statistics both converge almost surely to $c_0$. The realized reserve-truncated rent is bounded above by their gap, so dominated convergence yields
\begin{equation}
G_N^R\to v-c_0,
\qquad
R_N^R\to0.
\label{eq:reserve-limits}
\end{equation}
Appendix \ref{app:robustness} gives the complete argument. Hence if $A(v-c_0)>\Delta$, sufficiently large supplier markets again imply coordinated differentiation and unique decentralized duplication. We therefore interpret guaranteed trade as the condition supporting the strong monotone comparative static for every increment in $N$, not as a condition required for the long-run reversal itself.

\subsection{First-price procurement}

Second-price procurement is used in the baseline to make incidence transparent. Under standard symmetric independent-private-cost, risk-neutral, efficient-allocation conditions, revenue equivalence implies that first-price procurement generates the same expected payment and allocation objects relevant to the induced government game. In particular, expected buyer surplus and aggregate supplier rent are pinned down by the same order-statistic expectations. The upper-stage expected-payoff matrix is therefore unchanged under these standard conditions. This robustness does not extend automatically to risk aversion, asymmetric suppliers, common values, collusion, or endogenous participation.

\subsection{Why entry remains exogenous}

Endogenous supplier entry is intentionally outside the model. \citet{LiZheng2009} show that even in independent-private-value procurement, an increase in potential bidders can raise expected procurement cost because an entry effect may dominate the direct competition effect. Adding entry would require participation costs, an additional equilibrium, and a new welfare margin. The present theorem instead conditions on a supplier-market structure and asks how that market structure changes decentralized policy coordination. Endogenous entry is therefore a separate extension rather than a missing robustness check.

\section{Discussion and Empirical Implications}
\label{sec:discussion}

The model is stylized, but it generates several empirical predictions that distinguish the incidence channel from a conventional procurement-efficiency effect.

First, under the baseline incidence and DRHR conditions, greater supplier competition should weaken a jurisdiction's willingness to accept a supplier-side policy role when locally captured supplier rents are economically or politically salient. The relevant outcome is not whether procurement becomes cheaper conditional on a project. The prediction concerns whether the complementary activity is selected at all. Empirically, one would therefore look for changes in regional specialization, local-content support, supplier-development programs, or willingness to host complementary production stages as supplier concentration changes.

Second, the effect should be heterogeneous in local incidence. Regions whose governments internalize a larger share or weight of supplier-side surplus should have a larger effective $\rho$. They should therefore display a stronger negative relationship between supplier competition and willingness to accept the $D$ role. Conversely, when less supplier-side surplus enters the regional government's objective, $\rho$ is smaller and the rent-incidence mechanism is weaker in levels. If the jurisdiction also captures buyer-side gains---represented by $\sigma>0$ in Section \ref{sec:robustness}---the competition effect on local incentives can weaken or reverse.

Third, demand for intergovernmental coordination can increase precisely when procurement performance improves. In the model, a larger $N$ reduces expected production cost and improves total real surplus conditional on differentiated production. At the same time, it expands the parameter region in which decentralized governments duplicate $U$. This suggests an empirical distinction between \emph{within-project efficiency} and \emph{across-region policy coordination}. Better procurement outcomes need not imply better decentralized industrial-policy composition.

These predictions relate to, but differ from, the literature on procurement preferences. Bid preferences and set-asides deliberately alter the rules faced by favored firms \citep{Marion2007,AtheyCoeyLevin2013}. Our mechanism operates even when the procurement mechanism remains neutral: the distribution of supplier rents changes because the supplier market is more competitive. Nor is the mechanism one of local governments bidding against each other for firms, as in \citet{Slattery2025}. Here governments choose policy roles, and competition occurs among firms in the continuation market.

The theory also clarifies the appropriate policy interpretation. The model does not imply that governments should restrict supplier competition to preserve local rents. Doing so would sacrifice procurement efficiency and may introduce familiar favoritism and entry distortions. The result instead identifies a coordination problem: if competition erodes the local return to a socially valuable complementary role, the natural response may involve intergovernmental coordination, transfers, or institutional arrangements that reallocate some of the project surplus. Those instruments are deliberately not modeled here. Introducing them would change the policy mechanism and is left for future work.

Finally, the binary role choice should not be read as a literal claim that real regional policy is all-or-nothing. Its purpose is to isolate the equilibrium-regime effect. A continuous portfolio version would require specifying how procurement surplus scales with partial supplier-side policy and would move the paper toward the companion continuous-composition model. The present binary structure makes the market-structure comparative static and the disappearance of differentiated Nash equilibria exact.

\section{Conclusion}
\label{sec:conclusion}

Supplier competition improves procurement by lowering expected production cost. This paper shows that the same competition can worsen a different margin: decentralized coordination over which region undertakes a complementary industrial-policy role.

The mechanism relies on incidence. More suppliers raise expected real surplus from a complementary project but, under standard DRHR procurement conditions, reduce winning supplier information rent. A coordinated policymaker responds to the first effect, while a supplier-side regional government responds to the locally captured second effect. Consequently, competition lowers the coordinated threshold for regional differentiation but raises the decentralized threshold. The parameter region with unique inefficient duplication expands.

The stronger result is an equilibrium reversal. If differentiated production remains socially valuable when the supplier market becomes very competitive, then sufficiently strong supplier competition eliminates differentiated Nash equilibria and leaves duplicated local priorities as the unique decentralized outcome. The market becomes more efficient conditional on production while decentralized policy selection moves away from the coordinated allocation.

The auction-theoretic rent-dissipation result is not new, and neither is the generic observation that incomplete local capture can distort decentralized decisions. The contribution is their interaction across strategic layers: a familiar competition effect in the supplier market can change the equilibrium set of regional industrial policy. This mechanism suggests that procurement reform and industrial-policy coordination should be evaluated jointly when complementary production roles span jurisdictions.


\appendix
\section{Proofs}
\label{app:proofs}

\subsection{Proof of Proposition \ref{prop:nash}}

Fix $N$. If the other government chooses $D$, choosing $U$ yields $\Delta+A B_N>0$ whereas choosing $D$ yields zero, so $U$ is a strict best response to $D$.

If the other government chooses $U$, choosing $U$ yields $\Delta$ and choosing $D$ yields $A\rho R_N$. Therefore the best response to $U$ is $U$ when $A\rho R_N<\Delta$, both actions when equality holds, and $D$ when $A\rho R_N>\Delta$.

At the boundary $A\rho R_N=\Delta$, let $p$ be the probability that the opponent chooses $U$. Then
\begin{align}
\pi(U;p)-\pi(D;p)
&=\Delta+(1-p)A B_N-pA\rho R_N\\
&=(1-p)(\Delta+A B_N).
\end{align}
Hence $U$ is the unique best response for every $p<1$, whereas both $U$ and $D$ are best responses when $p=1$. Mutual best response therefore requires $p_1=1$ or $p_2=1$, which gives \eqref{eq:boundary-nash}. This also confirms that the three pure profiles listed in Proposition \ref{prop:nash} are the pure members of a larger boundary correspondence.

For $A\rho R_N>\Delta$, let $p$ be the probability that the opponent chooses $U$. The expected payoff from $U$ is
\begin{equation}
\pi(U;p)=p\Delta+(1-p)(\Delta+A B_N)
=\Delta+(1-p)A B_N,
\end{equation}
and the expected payoff from $D$ is
\begin{equation}
\pi(D;p)=pA\rho R_N.
\end{equation}
Indifference requires
\begin{equation}
\Delta+A B_N=pA(B_N+\rho R_N),
\end{equation}
which yields \eqref{eq:mixed}. Positivity is immediate, and $p_N^*<1$ iff $\Delta<A\rho R_N$. Hence this is the unique completely mixed equilibrium in the high-$A$ regime. Any semi-mixed equilibrium would require one player to use both actions while the other uses a pure action; but against $D$, $U$ is a strict best response, and against $U$, $D$ is a strict best response in this regime. Thus no such semi-mixed equilibrium exists. \qed

\subsection{Proof of Proposition \ref{prop:planner}}

The only positive candidate aggregate values are $2\Delta$ at $(U,U)$ and $\Delta+A G_N$ at either differentiated profile. Their difference is
\begin{equation}
V(U,D)-V(U,U)=A G_N-\Delta.
\end{equation}
The sign changes at $A=\Delta/G_N=A_N^P$. The profile $(D,D)$ yields zero and is never strictly optimal because $\Delta>0$. The result follows. \qed

\subsection{Proof of Theorem \ref{thm:divergence}}

Because $C_{1:N+1}\le C_{1:N}$ almost surely under the natural coupling, with strict inequality on an event of positive probability for a nondegenerate continuous distribution,
\begin{equation}
m_{N+1}<m_N.
\end{equation}
Hence $G_{N+1}=v-m_{N+1}>v-m_N=G_N$, and therefore
\begin{equation}
A_{N+1}^P=\frac{\Delta}{G_{N+1}}
<\frac{\Delta}{G_N}=A_N^P.
\end{equation}
Under DRHR, the procurement result in \citet{Li2005} and \citet{XuLi2019} gives $R_{N+1}<R_N$. Since $\rho>0$,
\begin{equation}
A_{N+1}^N=\frac{\Delta}{\rho R_{N+1}}
>\frac{\Delta}{\rho R_N}=A_N^N.
\end{equation}
\qed

\subsection{Proof of Corollary \ref{cor:wedge-expansion}}

First note that $G_N=B_N+R_N>\rho R_N$ because $B_N>0$ and $\rho\le1$. Hence $A_N^P<A_N^N$ and $\mathcal W_N$ is nonempty. Theorem \ref{thm:divergence} lowers its left endpoint and raises its right endpoint strictly when $N$ increases by one. Therefore $\mathcal W_N\subsetneq\mathcal W_{N+1}$. \qed

\subsection{Proof of Theorem \ref{thm:reversal}}

Under the compact support and positive density, $C_{1:N}\to c_0$ almost surely. Bounded convergence gives
\begin{equation}
m_N\to c_0,
\end{equation}
so $G_N\to v-c_0$. If $A(v-c_0)>\Delta$, there exists $N_P$ such that
\begin{equation}
A G_N>\Delta
\end{equation}
for all $N\ge N_P$. Proposition \ref{prop:planner} then implies that the coordinator strictly differentiates.

Under DRHR, $R_N\to0$ \citep{XuLi2019}. Hence there exists $N_N$ such that
\begin{equation}
A\rho R_N<\Delta
\end{equation}
for all $N\ge N_N$. Proposition \ref{prop:nash} then implies that $(U,U)$ is the unique Nash equilibrium. Taking $\bar N=\max\{N_P,N_N\}$ proves the result. \qed

\subsection{Proof of Corollary \ref{cor:destruction}}

At $N_0$, $A\rho R_{N_0}>\Delta$ implies differentiated pure Nash equilibria by Proposition \ref{prop:nash}. Because $R_N\to0$, there exists a finite $N_1>N_0$ such that $A\rho R_N<\Delta$ for all sufficiently large $N\ge N_1$, so only $(U,U)$ remains. If the planner differentiates at $N_0$, then $A G_{N_0}>\Delta$. Since $G_N$ is strictly increasing in $N$, $A G_N>\Delta$ for all $N>N_0$, so planner differentiation persists. \qed

\section{Additional Robustness Derivations}
\label{app:robustness}

\subsection{General incidence}

Let the $D$ government's local benefit be $L_N=\rho R_N+\sigma B_N$. Define $d_m=m_N-m_{N+1}>0$ and $d_R=R_N-R_{N+1}>0$. Since $s_N=m_N+R_N$,
\begin{equation}
B_{N+1}-B_N=s_N-s_{N+1}=d_m+d_R.
\end{equation}
Therefore
\begin{equation}
L_{N+1}-L_N=-\rho d_R+\sigma(d_m+d_R).
\end{equation}
Thus $L_{N+1}<L_N$ iff
\begin{equation}
\rho>\sigma\left(1+\frac{d_m}{d_R}\right),
\end{equation}
which is \eqref{eq:incidence-condition}.

\subsection{Uniform order statistics}

If $C_k\sim U[0,1]$, standard order-statistic formulas give
\begin{equation}
\E[C_{1:N}]=\frac{1}{N+1},
\qquad
\E[C_{2:N}]=\frac{2}{N+1}.
\end{equation}
Hence $R_N=1/(N+1)$, $B_N=v-2/(N+1)$, and $G_N=v-1/(N+1)$, yielding \eqref{eq:uniform-thresholds}.

\subsection{Reserve-value limits}

When trade occurs only if $C_{1:N}<v$ with $v>c_0$, define $G_N^R$ and $R_N^R$ as in \eqref{eq:reserve-G}--\eqref{eq:reserve-R}. Under the baseline compact support, continuity, and strictly positive density, for each fixed $k$ the $k$th order statistic converges almost surely to the lower endpoint; in particular,
\begin{equation}
C_{1:N}\to c_0,
\qquad
C_{2:N}\to c_0
\quad\text{a.s.}
\end{equation}
Therefore $(v-C_{1:N})_+\to v-c_0$ almost surely. This sequence is bounded on the compact support, so bounded convergence gives $G_N^R\to v-c_0$.

For supplier rent, define the realized reserve-truncated rent
\begin{equation}
r_N^R\equiv
(\min\{C_{2:N},v\}-C_{1:N})
\1\{C_{1:N}<v\}.
\end{equation}
On the active event it is nonnegative, and for every realization
\begin{equation}
0\le r_N^R\le C_{2:N}-C_{1:N}\le \bar c-c_0.
\end{equation}
Because both order statistics converge almost surely to $c_0$, their gap converges almost surely to zero, and hence $r_N^R\to0$ almost surely. Dominated convergence then yields
\begin{equation}
R_N^R=\E[r_N^R]\to0.
\end{equation}
This proves \eqref{eq:reserve-limits} under the stated assumptions and yields the eventual reversal whenever $A(v-c_0)>\Delta$. We do not claim one-step monotonicity of $R_N^R$.

\section{Computational Verification}
\label{app:verification}

The analytical proofs above are the basis for every theorem. For reproducibility, a Python/SymPy regression script independently checks the algebra of the mixed-equilibrium condition, the boundary best-response identity, the Uniform closed-form thresholds, and the general-incidence inequality. A separate Python script reproduces Figure \ref{fig:uniform-regimes}. These computations are verification aids and are not used as substitutes for the analytical proofs.

\section{Connection to the Companion Portfolio Model}
\label{app:companion}

At a fixed $N$, compare $(U,U)$ with a unilateral switch to $(U,D)$. The coordinated gain is
\begin{equation}
A G_N-\Delta,
\end{equation}
and the $D$ government's gain is
\begin{equation}
A\rho R_N-\Delta.
\end{equation}
Thus a reduced-form binary switching model with total complementary value $G^{\mathrm{red}}=A G_N$ and local capture share
\begin{equation}
\lambda_N=\frac{\rho R_N}{G_N}
\end{equation}
reproduces the fixed-$N$ threshold comparison. This is why the fixed-$N$ wedge is not treated as a separate contribution. The present paper's new comparative static is that the procurement market endogenously makes $G_N$ rise and $\lambda_N$ fall with supplier competition and thereby changes the exact regional-policy equilibrium set.


\section*{Data availability}
No data were used for the research described in this article. The algebraic verification scripts and figure-generation code are maintained with the reproducibility materials for this manuscript.

\section*{Declaration of generative AI and AI-assisted technologies in the manuscript preparation process}
During the preparation of this work, the author used OpenAI ChatGPT to assist with manuscript drafting and editing, algebraic and proof checking, literature organization, and the generation and review of verification code. After using this tool, the author reviewed and edited the content as needed, independently checked the mathematical claims and references used in the manuscript, and takes full responsibility for the content of the publication.

\begin{thebibliography}{99}

\bibitem[Aghion et al.(2015)]{AghionEtAl2015}
Aghion, P., Cai, J., Dewatripont, M., Du, L., Harrison, A., and Legros, P. (2015).
Industrial policy and competition.
\emph{American Economic Journal: Macroeconomics}, 7(4), 1--32.

\bibitem[Athey et al.(2013)]{AtheyCoeyLevin2013}
Athey, S., Coey, D., and Levin, J. (2013).
Set-asides and subsidies in auctions.
\emph{American Economic Journal: Microeconomics}, 5(1), 1--27.

\bibitem[Chiappinelli et al.(2025)]{ChiappinelliGiuffridaSpagnolo2025}
Chiappinelli, O., Giuffrida, L. M., and Spagnolo, G. (2025).
Public procurement as an innovation policy: Where do we stand?
\emph{International Journal of Industrial Organization}, 100, 103157.

\bibitem[Eckersley et al.(2023)]{EckersleyEtAl2023}
Eckersley, P., Flynn, A., Lakoma, K., and Ferry, L. (2023).
Public procurement as a policy tool: The territorial dimension.
\emph{Regional Studies}, 57(10), 2087--2101.

\bibitem[Justman et al.(2005)]{JustmanThisseVanYpersele2005}
Justman, M., Thisse, J.-F., and van Ypersele, T. (2005).
Fiscal competition and regional differentiation.
\emph{Regional Science and Urban Economics}, 35(6), 848--861.

\bibitem[Li(2005)]{Li2005}
Li, X. (2005).
A note on expected rent in auction theory.
\emph{Operations Research Letters}, 33(5), 531--534.

\bibitem[Li and Zheng(2009)]{LiZheng2009}
Li, T. and Zheng, X. (2009).
Entry and competition effects in first-price auctions: Theory and evidence from procurement auctions.
\emph{Review of Economic Studies}, 76(4), 1397--1429.

\bibitem[Lin and Li(2026)]{LinLi2026}
Lin, C. and Li, Y. (2026).
A unified national market and local industrial policy competition.
\emph{Frontiers of Economics in China}, 21(1), 1--32.

\bibitem[Marion(2007)]{Marion2007}
Marion, J. (2007).
Are bid preferences benign? The effect of small business subsidies in highway procurement auctions.
\emph{Journal of Public Economics}, 91(7--8), 1591--1624.

\bibitem[Motta and Polo(2026)]{MottaPolo2026}
Motta, M. and Polo, M. (2026).
Supply chain disruption and precautionary industrial policy.
\emph{International Journal of Industrial Organization}, 104, 103239.

\bibitem[Macatangay(2016)]{Macatangay2016}
Macatangay, R. E. M. (2016).
Optimal local content requirement policies for extractive industries.
\emph{Resources Policy}, 50, 244--252.

\bibitem[OECD(2026)]{OECD2026}
OECD (2026).
\emph{Public Procurement, Trade and Industrial Policies: Supporting National Priorities}.
OECD Publishing, Paris.

\bibitem[Saito and Mitsuta(2012)]{SaitoMitsuta2012}
Saito, T. and Mitsuta, N. (2012).
Analyzing local government procurements in Japan: Review of regional requirements.
\emph{Studies in Regional Science}, 42(2), 363--376.

\bibitem[Slattery(2025)]{Slattery2025}
Slattery, C. (2025).
Bidding for firms: Subsidy competition in the United States.
\emph{Journal of Political Economy}, 133(8), 2563--2614.

\bibitem[Szech(2011)]{Szech2011}
Szech, N. (2011).
Optimal advertising of auctions.
\emph{Journal of Economic Theory}, 146(6), 2596--2607.

\bibitem[Vagstad(1995)]{Vagstad1995}
Vagstad, S. (1995).
Promoting fair competition in public procurement.
\emph{Journal of Public Economics}, 58(2), 283--307.

\bibitem[Van Biesebroeck and Zhang(2026)]{VanBiesebroeckZhang2026}
Van Biesebroeck, J. and Zhang, H. (2026).
Lessons on industrial policy from industrial organization: Introduction to the special issue.
\emph{International Journal of Industrial Organization}, 104, 103243.

\bibitem[Xu and Li(2019)]{XuLi2019}
Xu, M. and Li, D. Z. (2019).
Equilibrium competition, social welfare and corruption in procurement auctions.
\emph{Social Choice and Welfare}, 53, 443--465.

\bibitem[Anonymous(2026)]{Companion2026}
Anonymous (2026).
Industrial policy composition and regional value chains.
Companion manuscript.

\end{thebibliography}


\end{document}
